\documentclass[12pt]{article}
\usepackage[margin=1in]{geometry}
\usepackage{amsmath,amssymb,amsthm,mathtools}
\usepackage{booktabs,threeparttable,graphicx,setspace,enumitem,longtable,array}
\usepackage{xcolor}
\usepackage{etoolbox}
\usepackage{placeins}
\AtBeginEnvironment{tabular}{\footnotesize}
\usepackage[authoryear,round]{natbib}
\usepackage[colorlinks=true,linkcolor=blue!50!black,citecolor=blue!50!black,urlcolor=blue!50!black]{hyperref}
\usepackage{fancyhdr}
\doublespacing
\emergencystretch=2em
\pagestyle{fancy}\fancyhf{}
\fancyhead[L]{\small Niroula, Alkhushayni, Chen, Bhavsar, and Niroula}
\fancyhead[R]{\small From Human Markets to Machine Markets}
\fancyfoot[C]{\thepage}
\renewcommand{\headrulewidth}{0.3pt}
\setlength{\headheight}{14pt}

\newtheorem{theorem}{Theorem}[section]
\newtheorem{proposition}[theorem]{Proposition}
\newtheorem{lemma}[theorem]{Lemma}
\newtheorem{corollary}[theorem]{Corollary}
\theoremstyle{definition}
\newtheorem{definition}[theorem]{Definition}
\theoremstyle{remark}
\newtheorem{remark}[theorem]{Remark}
\newtheorem{example}[theorem]{Example}
\newcommand{\lab}[1]{\textsc{[#1]}}
\newcommand{\R}{\mathbb{R}}
\providecommand{\E}{\mathbb{E}}

\title{From Human Markets to Machine Markets:\\
Transient Amplification, an Exact Depth Condition,\\ and a Maximal Pre-Trade Guard}
\author{Pratik Niroula\thanks{Department of Computer Information Science, Minnesota State University, Mankato. \texttt{pratik.niroula@mnsu.edu}.}
\and Suboh Alkhushayni\thanks{Department of Computer Information Science, Minnesota State University, Mankato. \texttt{suboh.alkhushayni@mnsu.edu}.}
\and Yuhao Chen\thanks{Department of Accounting and Finance, Minnesota State University, Mankato. \texttt{yuhao.chen@mnsu.edu}.}
\and Mansi Bhavsar\thanks{Department of Computer Information Science, Minnesota State University, Mankato. \texttt{mansi.bhavsar@mnsu.edu}.}
\and Prashant Niroula\thanks{Faculty of Management, Tribhuvan University, Kirtipur, Nepal. \texttt{prashantniroula20@gmail.com}.}}
\date{Revised preprint, October 2026}

\begin{document}
\maketitle
\thispagestyle{empty}

\begin{abstract}
\noindent Eigenvalues and volatility do not bound transient losses before a market settles. We show that cumulative transient risk in a limit order book equals $S=\tilde\kappa/(m-q\tilde\kappa)$, that a book overshoots if and only if $\varphi\varepsilon>(1-a)\kappa$, and that a pre-trade guard on the summed impulse response is maximal and invariant to unobserved trader cohorts (including under carried inventory and interval-bounded hidden drift). Certifying $S$ requires $O(\sigma^2\epsilon^{-2}\log(1/\epsilon))$ experiments regardless of cohort spacing, yielding a minimax-optimal guard. Across 25 pre-registered primary tests, 15 pass: in thin Binance perpetuals, six-hour estimates forecast next-hour adverse excursions out of sample across five tests ($\Delta R^2=+0.85$ to $+4.20$ percentage points; $+2.79$ [$+1.43, +4.94$] beyond HAR/EWMA volatility and $+0.56$ [$+0.27, +1.16$] beyond VPIN), replicate in thin US equities ($+0.50$ [$+0.27, +0.78$]), and lower tail loss by $8.1\%$ ($[-16.1\%, -1.7\%]$) and $22.6\%$ ($[-39.2\%, -2.0\%]$) against static limits.
\end{abstract}

\noindent\textbf{Keywords:} transient amplification, non-normal dynamics, limit order book, market resilience, liquidity externality, pre-trade risk control, pre-registration, replication.\\
\textbf{JEL:} G01, G10, G14, G18, C32, C58.

\paragraph{Reproducibility.} Every number in the paper is produced by a script or notebook cell that writes a result file from public data or from a closed form; Appendix~\ref{app:repro} maps each table, figure and number to its source. Pre-registered specifications, their SHA-256 hashes and test-set unlock counters are in the replication package. Every simulation number was re-run for this version.

\newpage
\setcounter{page}{1}
\tableofcontents
\newpage

\section{Introduction}

On October 10, 2025, crypto derivatives markets liquidated an estimated \$19 billion of leveraged positions in a single day, the largest such event on record \citep{shares2025}; on May 6, 2010, US equity prices fell and recovered by several percent within minutes \citep{kirilenko2017}. In both episodes the losses that mattered were transient: positions were closed out on the path, before prices settled. This paper is about that path risk in markets where algorithms trade with each other.

Algorithmic markets are closed-loop systems: trading moves prices, prices move algorithms, and algorithms trade. A common shortcut equates the stability of such a system with its eigenvalues lying inside the unit circle. That is the wrong object for a leveraged intermediary, whose survival depends on the path of losses before the system settles. A market can be spectrally stable and still push a dealer through its maintenance margin within seconds.

This paper asks three questions. When does a limit order book amplify a flow shock, so that prices keep moving after the shock has passed? How can a dealer or a venue control transient loss without identifying who is trading, which the companion paper \citep{niroula2026} shows is expensive or impossible from public data? And when the traders cannot be identified, what information certifies a safe trading decision, what does that information cost, and how much safe trading must be given up when it is missing? We give exact answers in a linear book model and then confront them with data, reporting the results as they came out.

\paragraph{1. Spectral stability does not bound transient loss.} The exact safe shock radius of a stable linear market is the loss budget divided by its peak transient gain, and that gain can be arbitrarily large at any fixed spectral radius (Theorem~\ref{thm:radius}). Markets with identical eigenvalues, one-step volatility and stationary price variance can have arbitrarily large worst-case transient gain, as the companion paper shows \citep{niroula2026}. Supervision through eigenvalues or volatility alone is therefore insufficient in principle.

\paragraph{2. When a book amplifies.} We distinguish peak overshoot, a response that rises above its direct impact, from cumulative amplification, a summed response far above it. When order flow depletes one side of the book and the book refills at rate $r$, the price response overshoots its direct impact if and only if $\varphi\varepsilon>(1-a)\kappa$ (Theorem~\ref{thm:peak}\ref{thm:t6}). With momentum traders the one-period margin $C_3=\varphi\varepsilon+\kappa^2\eta q-(1-a)\kappa$ is sufficient (Theorem~\ref{thm:peak}\ref{prop:t6m}) but no one-period condition is exact (Remark~\ref{rem:counter}); momentum alone cannot make a stable book overshoot (Theorem~\ref{thm:peak}\ref{thm:nomom}), and an explicit condition rules overshoot out (Theorem~\ref{thm:peak}\ref{thm:noover}). Cumulative risk has a closed form, $S=\tilde\kappa/(m-q\tilde\kappa)$: long-run impact divided by reversion net of momentum (Theorem~\ref{thm:sum}). It needs no depletion, and it is the form our evidence supports: crypto books show no peak overshoot, but thin books revert slowly, so their summed response is large.

\paragraph{3. The maximal safe guard.} For bounded shocks, worst-case loss over a horizon is inventory times the shock bound times the summed absolute impulse response, the classical $\ell_1$ gain of control theory. Under a declared risk contract (fixed inventory, per-period shock bound, mark-to-market loss), a pre-trade guard on this summed gain is necessary and sufficient for staying within the loss budget, so it admits the largest safe set of actions (Lemma~\ref{thm:guard}). A guard on the peak response, the quantity closest to the eigenvalue intuition, admits unsafe actions. A Gr\"onwall bound carries the guard to an estimated operator (Theorem~\ref{thm:robust}), a contraction-block certificate makes it computable on a finite horizon (Lemma~\ref{lem:tail}), and its inputs are identified from prices without identifying the algorithms (Proposition~\ref{prop:invariance} and Theorem~\ref{thm:certcost}). We embed the guard in a dealer's balance sheet with margin, hedging and fill accounting (Theorem~\ref{thm:forward}) and show that strategic re-optimization by filtering algorithms reinforces it (Proposition~\ref{prop:lucas}).

\paragraph{4. Information and safety.} The information that certifies a safe decision is not the information that identifies the traders. Two markets can agree exactly on their early price responses and differ threefold in risk (Example~\ref{ex:witness}), and any supervisor who must rarely admit an unsafe level then gives up a computable amount of safe capacity (Proposition~\ref{prop:twoworld}). A finite window of responses certifies the summed response to an interval computed exactly from the Markov--Krein extremal measures, which shrinks geometrically with the window under a bound on memory length and is unbounded without one (Theorem~\ref{thm:window}). With noisy reset experiments the summed response is certified at the parametric rate, $N=O(\sigma^2\epsilon^{-2}\log(1/\epsilon))$, whatever the number and spacing of the traders' memories, while identifying them costs $(\Delta\eta)^{-(4K-2)}$ or more (Theorem~\ref{thm:certcost}); the companion paper's hardest pairs have identical summed responses. A guard built on the certificate has a closed form, admits an interval of levels, flags infeasible states, and is minimax-optimal in lost capacity up to a logarithm (Theorem~\ref{thm:robustguard}). When admission also changes the book, the safe levels need not form an interval, so certifying the largest one no longer certifies those below it; we certify the whole safe set by interval branch and bound (Proposition~\ref{prop:venue}, Theorem~\ref{thm:venuebb}). The window bounds are proved rather than computed (Lemma~\ref{lem:certnum}), and certificates that allow correlated noise keep their level where the simple one fails (Proposition~\ref{prop:hotelling}).

\paragraph{5. The value of disclosure.} Disclosure that rules out the extremal compatible economy raises safe capacity, and its value grows as memory lengthens, approaching the full oracle capacity; disclosure of a noisy proxy can be worth less than nothing (Theorem~\ref{thm:disclosure}, Table~\ref{tab:discsim}).

\paragraph{6. Who controls the risk.} Reversion is a public good. In equilibrium dealers under-provide it, momentum traders over-provide feedback, fewer dealers mean more transient risk, and the rebate that restores the first best is $(1-\theta/N)\Lambda S^2/\tilde\kappa$, a function of the guard's own statistic (Proposition~\ref{prop:equil}). When dealers also choose depth, fewer dealers raise risk through larger impact while the unit-free multiple $S/\tilde\kappa$ stays fixed, the pattern we see across Binance perpetuals.

\paragraph{7. Evidence.} On Binance BTCUSDT futures the amplification statistic is negative on 93.5\% of days, depth does not carry an overshoot, and the statistic fails a pre-registered forecasting test (Section~\ref{sec:depth}); unit-free spot amplification is indistinguishable from one (Section~\ref{sec:spot}). Across 30 perpetuals, Binance's depth bands, which are re-centred on the current mid, show taker flow replenishing the side it hits; on a fixed price grid in Bybit's full order books, the same flow depletes it, so the replenishment finding is an artefact of the band construction (Section~\ref{sec:fixedgrid}). The amplifying regime is more frequent in thin books (Section~\ref{sec:thin}). There, six-hour estimates of the amplification margin forecast the size of the next hour's adverse excursion out of sample in four pre-registered tests: on the original ten symbols ($+1.2$ points of $R^2$), on ten new symbols ($+0.85$), on fresh days with price-impact controls ($+1.6$), and on further fresh days beyond a full volatility-forecasting baseline ($+2.8$, with better tail classification). The summed response $S$, the guard's own statistic, carries the same information ($+2.1$ points beyond the volatility baseline), the gain survives adding VPIN to the baseline ($+2.0$ points post hoc, $+0.56$ in a pre-registered test on new symbols), unit-free versions work, and the content comes mostly from impact-weighted reversion, as Theorem~\ref{thm:sum} implies. Within a volatility quintile, hours in the top quintile of the statistic have next-hour excursions about 26 basis points larger than hours in the bottom quintile, enough to roughly double the hourly liquidation probability of a $20\times$ position. The effect is much weaker in mid-depth books and absent in BTCUSDT. On Bybit it fails with trade data alone; with order-book depth it meets a pre-registered criterion, but the gain is about one eightieth of Binance's (Table~\ref{tab:bybitob}). In a regulated market, 23 thin US stocks with consolidated trades, the depth-free statistic passes a fifth pre-registered forecasting test ($+0.50$ points of $R^2$ beyond the same volatility baseline), again through impact-weighted reversion; unlike on Binance, mid-size stocks show a gain too (Table~\ref{tab:us}). As a position limit, a six-hour estimate of the response adds nothing to volatility, but once the operator is certified from the previous three days, scaling a volatility limit by the gain multiple lowers tail loss per unit of position by 2.4\% in a fifth pre-registered test, a narrow pass; limits built on the estimated response alone remain worse than volatility (Section~\ref{sec:certguard}).

Simulations with a known truth confirm the coverage, attainment and optimality claims and show where the certificate fails: when memory is longer than declared, and when feedback strengthens after certification (Section~\ref{sec:simtests}). On days that no earlier part used, two pre-registered primaries pass: the forecasting model, frozen at its 2023--2024 estimates, adds $4.2$ points of $R^2$ beyond the volatility baseline, and the uncertainty-aware guard lowers tail loss per unit of position by 8\% against a static limit at 91\% of its activity, doing much better than depth or spectral-radius throttles; it does not beat a volatility cap, and its bootstrap upper bound does no better than the point estimate (Section~\ref{sec:p15}). On thirty perpetuals that no earlier part used, of which a seeded rule kept five thin ones, three pre-registered primaries pass: $\hat C_3$ adds $0.73$ points of $R^2$ beyond the volatility baseline and $0.56$ beyond the baseline plus VPIN, turning the post-hoc VPIN comparison into a confirmatory one, and the guard lowers tail loss by 23\% against a static limit, with a wide interval (Section~\ref{sec:p16}). A pre-registered test on a venue no earlier part used, Coinbase spot, fails: in the six thin pairs that trade often enough, the statistic adds nothing (Section~\ref{sec:p17}). Post hoc, a placebo that keeps the flow's persistence but breaks its alignment with prices adds nothing on Binance, where the real statistic adds $2.7$ points, so the content is the alignment rather than order splitting; and the guard has lower tail loss than certified peak-response, margin-style and participation rules on the Part 15 sample, but not than a six-hour peak guard (Section~\ref{sec:p18}). Three further analyses were specified in an earlier draft before their data were downloaded but were not hash-registered (Section~\ref{sec:p1921}): on fresh days for 19 thin coins, the predicted interaction of $\hat C_3$ with open-interest-to-depth leverage fails ($-0.18$~bp, CI $[-3.02,+3.73]$) while the spot gain is $15.2\%$ of the perpetual gain ($+0.99$ vs.\ $+6.50$ points, ratio CI $[-0.015,+0.308]$; Section~\ref{sec:p19}); in Binance's BTCUSDT spot zero-fee experiment, the primary reversion and gain-multiple hypotheses are not supported while impact falls $18\%$ in the two spot events ($t=-3.86$, exploratory, single treated pair; Section~\ref{sec:p20}); and in dollars, the guard alone lowers hourly expected shortfall on \$1M of average position by \$3{,}045 and \$10{,}518 against a static limit on the untouched days and symbols, while its gain when combined with a volatility cap (\$398 and \$2{,}557) is not statistically significant (Section~\ref{sec:p21}). Finally, four pre-registered follow-up tests on untouched days (Sections~\ref{sec:p22}--\ref{sec:p25}) test the mechanism and multi-hour pre-trade guard directly: on untouched 22nd-of-month days (Part 22), $\hat C_3$ strongly forecasts next-hour adverse excursions ($\Delta R^2=+6.52$ points, $t=12.21$), but neither forecasts five-minute open-interest contractions ($\Delta R^2=-0.80$ points) nor is mediated by them (mediation share $-6.8\%$, both primaries fail), whereas on 42 untouched 2nd-and-9th-of-month days (Part 25, $13{,}532$ test symbol-hours), $\hat C_3$ strongly forecasts next-hour 15-minute checkpoint price-impact persistence ($\Delta R^2_{\rm pers}=+3.71$ points, $t=14.84$, Primary 1 passes; $+3.94$ points on excursions, $t=15.68$) and controlling for 15-minute persistence mediates $97.6\%$ of $\hat C_3$'s excursion $\Delta R^2$ ($[96.6\%,98.3\%]$; $98.6\%$ for the 60-minute $L_1$ integral, though the single five-minute bar jump placebo also exceeds $50\%$ so the joint Primary 2 criterion is recorded as a failure). And for non-resetting 18-hour leveraged positions, while unbounded unit-rescaled matched-breach capacity gains have wide intervals that cross zero ($+\$332{,}848$, CI $[-\$20{,}390,+\$1{,}180{,}808]$ in Part 23; $+\$172{,}416$, CI $[-\$219{,}311,+\$578{,}090]$ in Part 24), enforcing the SEC Rule 15c3-5 capped pre-trade throttle $q_t\in(0,1]$ on a $\$1\text{M}$ $20\times$ facility across 42 untouched test dates (Part 24) reduces the forced-liquidation rate by $-7.41$ percentage points beyond a volatility cap (CI $[-11.45,-4.15]$, and $-4.09$ percentage points per $\$1\text{M}$ deployed, CI $[-8.47,-1.03]$, Primary 1 passes) and cuts 18-hour daily $\mathrm{ES}_{95}$ dollar tail loss by $-\$26{,}429$ (CI $[-\$53{,}146,-\$6{,}254]$, Primary 2 passes).

\paragraph{What we do not claim.} The information-safety theorems concern declared observation and risk contracts (controlled reset experiments, exact windows, tagged or audited data); they are not finite-sample guarantees for observational order-book records, for which we report a bootstrap diagnostic and its out-of-sample performance. Our cohorts are any traders whose demand follows a linear filter of past prices, human or automated; nothing in the evidence attributes an effect to algorithms rather than to people. We do not claim that crypto books amplify shocks through depth: flow depletes the side it hits on a fixed price grid, but we observe no peak overshoot, and the depth term carries no forecasting content. The thin-book forecasting result is reduced-form, and its size depends on the venue: largest in thin Binance books, about a fifth as large in thin US stocks, negligible on Bybit and absent on Coinbase; it improves tail classification only against the full volatility baseline. We do not claim that the guard replaces volatility in practical position sizing: its theorems are statements about worst-case loss under a certified operator, an estimated six-hour response adds nothing to volatility, and a three-day certified estimate improves a volatility limit only modestly (2.4\%, a narrow pass) while the guard alone does worse. The October 10, 2025 episode is reported descriptively and supports no inference (Section~\ref{sec:oct10}). One test cell of the first thin-book test was executed twice with unchanged code and specification (Section~\ref{sec:x6}). The equilibrium is a stylized model; thinner books do carry more cumulative risk per dollar of flow, as it predicts, but through larger impact rather than through the weaker reversion it names.

\begin{table}[p]\centering\footnotesize
\begin{threeparttable}
\caption{Map of results}\label{tab:map}
\begin{tabular}{p{0.34\textwidth}p{0.38\textwidth}p{0.16\textwidth}}
\toprule
Question & Answer & Where\\
\midrule
Does $\rho(J)<1$ bound transient loss? & No; safe radius $B/M_C$ can vanish at fixed $\rho$ & Thm~\ref{thm:radius}, Companion paper\\
Can volatility monitors help? & No; fixed variance, unbounded gain & Companion paper\\
When does a book amplify? & Iff $\varphi\varepsilon>(1-a)\kappa$; with momentum $C_3>0$ suffices, no exact one-period rule & Thm~\ref{thm:peak}\ref{thm:t6}, Thm~\ref{thm:peak}\ref{prop:t6m}, Rem~\ref{rem:counter}\\
Can momentum alone amplify? & Not in a stable book & Thm~\ref{thm:peak}\ref{thm:nomom}\\
How large is cumulative risk? & $S=\tilde\kappa/(m-q\tilde\kappa)$ & Thm~\ref{thm:sum}\\
Which guard is safe and largest? & Summed absolute gain; peak gain is unsafe & Lem~\ref{thm:guard}\\
Robust to estimation error? & Yes, $M/(1-\epsilon S)$ & Thm~\ref{thm:robust}\\
Needs identification of algorithms? & No; inputs are invariant to cohort splits & Prop~\ref{prop:invariance}, Thm~\ref{thm:certcost}\\
Do learners undo the guard? & No; gains fall when admission falls, and every finite-horizon summed response rises with the gains in the cohort economy & Prop~\ref{prop:lucas}\\
Is reversion efficiently supplied? & No; under-provided, and the fix is priced by $S$ & Prop~\ref{prop:equil}\\
Is BTCUSDT on the amplifying side? & No: $\hat C_3<0$ on 93.5\% of days & Table~\ref{tab:d3}\\
Does depth carry an overshoot? & No; all overshoots negative & Figure~\ref{fig:med}\\
Does $\hat C_3$ forecast tail risk in BTCUSDT? & No (pre-registered) & Table~\ref{tab:d8}\\
Is spot amplification above one? & No; unit-free gain $\approx1.0000$ & Section~\ref{sec:spot}\\
Does flow deplete thin books? & Yes on a fixed price grid; the apparent replenishment is a band-recentring artefact & Tables~\ref{tab:x8}, \ref{tab:fixedgrid}\\
Does $\hat C_3$ forecast tail size in thin books? & Yes, in four pre-registered tests on Binance & Tables~\ref{tab:x6}, \ref{tab:y6}, \ref{tab:z6}, \ref{tab:z9}\\
Is it just volatility? & No: $+2.8$ pts beyond HAR, EWMA and hour effects; $S$ itself $+2.1$ & Table~\ref{tab:z9}\\
Is it just VPIN? & No: $+2.0$ pts beyond the baseline plus VPIN (post hoc); $+0.56$ on new symbols (pre-registered) & Section~\ref{sec:z9}, Table~\ref{tab:p16}\\
Is it tick size, impact or scale? & No; survives controls, unit-free versions work & Tables~\ref{tab:x6r}, \ref{tab:y8}, \ref{tab:z6}, \ref{tab:z9}\\
Does it hold on a second venue? & Bybit: absent with trades; with depth, significant but negligible & Section~\ref{sec:bybit}, Table~\ref{tab:bybitob}\\
In a regulated market? & Thin US stocks: yes, $+0.50$ pts (pre-registered); mid-size stocks too & Table~\ref{tab:us}\\
Does the estimated response size positions better? & Six-hour: no; certified three-day: $-2.4\%$ tail loss, narrow pass; alone: worse & Tables~\ref{tab:backtest}, \ref{tab:certguard}\\
How large is the effect? & top vs bottom quintile $\approx26$ bp at fixed volatility & Section~\ref{sec:threats}\\
Can equal early data hide different risk? & Yes, threefold; a computable capacity loss follows & Ex.~\ref{ex:witness}, Prop.~\ref{prop:twoworld}\\
What does a finite window certify? & An exact Markov--Krein interval; nothing without a memory bound & Thm~\ref{thm:window}\\
Certify risk vs identify traders? & $\sigma^2\epsilon^{-2}\log(1/\epsilon)$ vs $(\Delta\eta)^{-(4K-2)}$ & Thm~\ref{thm:certcost}, Table~\ref{tab:certvsid}\\
Optimal guard under uncertainty? & Closed form, interval, minimax capacity loss & Thm~\ref{thm:robustguard}\\
If admission changes the book? & Safe set can be two intervals; certified by branch and bound & Prop~\ref{prop:venue}, Thm~\ref{thm:venuebb}\\
If tap noise is correlated? & Simple certificate undercovers (75\%); $t$ and Hotelling certificates exact & Prop~\ref{prop:hotelling}\\
What is disclosure worth? & Most when memory is long; noisy proxies can hurt & Thm~\ref{thm:disclosure}, Table~\ref{tab:discsim}\\
On untouched days? & Yes: frozen model $+4.2$ pts; guard $-8\%$ vs static, not vs volatility (pre-registered) & Table~\ref{tab:p15}\\
On untouched symbols? & Yes, on five: $+0.73$ pts; $+0.56$ beyond VPIN; guard $-23\%$ vs static (pre-registered) & Table~\ref{tab:p16}\\
On a new venue? & No: Coinbase, $-0.09$ pts, fails (pre-registered) & Table~\ref{tab:p17}\\
Is it order-flow persistence? & No: a placebo with the same flow persistence adds nothing on Binance (post hoc) & Table~\ref{tab:p18}\\
Against peak, participation and margin rules? & Better than certified peak, margin and participation rules on Part 15; not than six-hour peak (post hoc) & Table~\ref{tab:p18}\\
Is leverage or OI liquidation the mechanism? & No: OI drops $-6.8\%$ mediation; multi-min persistence $+3.71$ pts ($97.6\%$ mediation) & Tables~\ref{tab:p19}, \ref{tab:p22}, \ref{tab:p25}\\
Do zero fees change reversion or impact? & Primary reversion/multiple null; spot impact $-18\%$ (exploratory, 1 treated pair) & Table~\ref{tab:p20}\\
What is the guard worth in dollars and capacity? & Capped 18h throttle: $-7.41$ pp liq.\ rate, $-\$26.4$k ES$_{95}$ vs vol.\ cap; unbounded $\overline{N}$ CI crosses 0 & Tables~\ref{tab:p21}, \ref{tab:p23}, \ref{tab:p24}\\
\bottomrule
\end{tabular}
\end{threeparttable}
\end{table}

\subsection{Related literature}

\paragraph{Closest work.} Three strands are closest. Models of transient impact and resilience \citep{obizhaeva2013,bouchaud2004,bouchaud2018} take the decay of impact as given; we derive when depletion, refill and momentum make a book overshoot, give cumulative impact in closed form, and ask who supplies the reversion that bounds it. Models of liquidity spirals \citep{brunnermeier2009,shleifer1992} explain how funding constraints amplify shocks; our amplification is a property of book dynamics and does not need funding constraints, and our guard is the pre-trade control that bounds the resulting loss. \citet{sornette2023} show that non-normal interactions create transient bubbles; we give a necessary and sufficient pre-trade control for such transients and test the conditions for them in order-book data.

\paragraph{Market microstructure and resilience.} Transient impact and resilience follow \citet{obizhaeva2013}, \citet{bouchaud2004} and \citet{bouchaud2018}; optimal execution with impact is in \citet{almgren2001} and \citet{cartea2015}. Depth imbalance as a short-horizon price predictor is documented by \citet{cont2014}. \citet{biais2005} survey the field; price impact as a measure of information and illiquidity goes back to \citet{kyle1985} and \citet{amihud2002}. Inventory-based dealer models go back to \citet{ho1981}, and \citet{avellaneda2008} give the standard high-frequency market-making model. Empirically, resiliency, the speed at which a book refills after a liquidity shock, is measured by \citet{biais1995}, \citet{degryse2005} and \citet{large2007}, who estimates it with a self-exciting point process for order events; \citet{foucault2005} show that it is set by the mix of patient and impatient traders and the tick size. Our reversion term $m$ is the reduced-form counterpart of resiliency, and $S=\tilde\kappa/(m-q\tilde\kappa)$ shows how it combines with impact and momentum: a book that is resilient in this sense can still carry a large cumulative response when momentum feedback absorbs most of its reversion. Our book model is a minimal two-state version of resilience dynamics, chosen so that the amplification condition is exact.

\paragraph{Feedback, fire sales and automated markets.} Liquidity spirals and fire sales are in \citet{shleifer1992} and \citet{brunnermeier2009}; automated-market crashes in \citet{kirilenko2017}; and the market-design case for discrete auctions in \citet{budish2015}. \citet{menkveld2013} documents high-frequency market makers as the new intermediaries, and algorithmic trading has been found to improve liquidity and price discovery in normal times \citep{hendershott2011,brogaard2014}; our question is what happens on the path after a shock. \citet{sornette2023} show that non-normal interactions can generate transient bubbles in unleveraged systems. We add an exact condition for when a book amplifies, a control that does not need the identity of the traders, and evidence on whether the condition holds in a large market.

\paragraph{Early warning and liquidity crashes.} The best-known early-warning metric for flow-driven crashes is VPIN \citep{easley2012}, whose forecasting content beyond volume and volatility is disputed \citep{andersen2014}. Equilibrium models explain crashes through endogenous liquidity withdrawal \citep{huang2009} and illiquidity contagion across assets \citep{cespa2014}. Our statistic differs in kind: it is a structural quantity, the summed response of the book to a unit of flow, with a closed form and a control-theoretic meaning, and we evaluate it out of sample against a full volatility-forecasting baseline. In a post-hoc comparison on the same test hours, $\hat C_3$ keeps about 73\% of its gain once VPIN is added to the baseline (Section~\ref{sec:z9}). Crypto markets are studied as a laboratory for microstructure and arbitrage by \citet{makarov2020}.

\paragraph{Non-normal dynamics and robust control.} Transient growth of stable non-normal systems and the Kreiss matrix theorem are classical \citep{trefethen2005,kreiss1962}; the sharp $en$ constant is due to \citet{spijker1991}. Standard $H_\infty$ control bounds frequency-domain gain under known coordinates. The time-domain worst case under bounded disturbances, the $\ell_1$ norm of the impulse response, is the object of $\ell_1$-optimal control \citep{vidyasagar1986,dahleh1987,dahleh1995}, and our guard's worst-case identity is that classical fact. What we add is its use as an admission rule on a dealer's balance sheet, its invariance to the unobservable cohort coordinates and identification from prices, and its link to who supplies reversion in equilibrium. Cross-impact under no dynamic arbitrage is studied by \citet{schneider2019}.

\paragraph{Regulation.} SEC Rule 15c3-5 requires broker-dealers with market access to enforce pre-trade controls on credit and capital \citep{sec2010}. Our guard is a candidate pre-trade control with a provable maximality property.

\paragraph{Forecast evaluation.} Our predictive tests use \citet{clark2007} and \citet{diebold1995} with \citet{newey1987} standard errors and block resampling in the spirit of \citet{kunsch1989}, the volatility baseline follows \citet{corsi2009} and \citet{riskmetrics1996}, local projections follow \citet{jorda2005}, and we correct for multiple testing as in \citet{holm1979}. Following the case for pre-registration in science \citep{nosek2018} and the concerns about multiple testing in financial economics \citep{harvey2016,harvey2017}, we pre-register specifications, count every evaluation of a held-out period, and report the failures with the successes.

\subsection{Outline}
Section~\ref{sec:model} sets out the closed-loop market. Section~\ref{sec:radius} shows why spectral stability is not solvency. Section~\ref{sec:t6} derives the amplification condition. Section~\ref{sec:guard} constructs and analyzes the guard, and Section~\ref{sec:equil} asks who supplies reversion in equilibrium. Sections~\ref{sec:depth} and~\ref{sec:spot} report the BTCUSDT evidence, and Section~\ref{sec:thin} the thin-book program, its replications and its failures. Section~\ref{sec:policy} discusses policy, and Section~\ref{sec:conclusion} concludes.
